

%% [TODO] 대괄호는 러닝 헤드(짧은 제목)다.
\title[Calibration collapse in time-split validation]{Beyond AUROC: Quantifying
Calibration Collapse in Time-Split Validation of Hepatotoxicity and
Cardiotoxicity Models}

%%=============================================================%%
%% GivenName	-> \fnm{Joergen W.}
%% Particle	-> \spfx{van der} -> surname prefix
%% FamilyName	-> \sur{Ploeg}
%% Suffix	-> \sfx{IV}
%% \author*[1,2]{\fnm{Joergen W.} \spfx{van der} \sur{Ploeg} 
%%  \sfx{IV}}\email{iauthor@gmail.com}
%%=============================================================%%

%% 단독 저자. 템플릿의 둘째·셋째 저자와 \equalcont 는 지웠다.
%% [TODO] 이름·이메일을 채울 것. ORCID 는 투고 시스템에 입력한다.
\author*[1]{\fnm{[이름]} \sur{[성]}}\email{[이메일]}

%% [TODO] 소속 한 곳.
\affil*[1]{\orgdiv{[학과]}, \orgname{[기관]}, \orgaddress{\city{[도시]}, \country{[국가]}}}

%%==================================%%
%% Sample for unstructured abstract %%
%%==================================%%

%%===========================================================%%
%% Journal of Cheminformatics: abstract must be UNSTRUCTURED, %%
%% 350 words maximum, and must contain a section titled      %%
%% "Scientific Contribution" of at most THREE sentences.     %%
%%===========================================================%%

\abstract{Models for drug-induced liver injury and cardiotoxicity are usually validated
on random splits and reported by the area under the receiver operating
characteristic curve. Deployment is
prospective: a model trained on approved drugs is applied to drugs approved
later. We asked what a random split hides, using two public FDA datasets (DILIrank,
796 compounds; DICTrank, 919) split at 2010, with compounds, features and
learner held fixed on both sides of every comparison.

Discrimination fell modestly, from 0.735 to 0.648 and from 0.697 to 0.620. A Murphy
decomposition of the Brier score attributed 86.1\% and 82.5\% of the loss to
reliability rather than resolution: the cost falls on the probabilities, not the
ranking. Because the positive rate also drops across the boundary, we built random
controls matched on held-out size and on one or both base rates; after matching
both, the temporal expected calibration error was still 1.37 and 1.21 times the
control.
Recalibrating inside the training era, as a prospective user must, recovered
8\% and 42\% of the distance to an oracle ceiling; split conformal prediction
missed its nominal 90\% coverage, at 0.737 and 0.824; and eight further
remedies, from prior-probability correction and nine quantification estimators
to density-ratio weighting and four additional calibrators, did not close the
gap. One intervention did.
Retraining as approvals accumulate and recalibrating at each refit lowered the
error by 0.166 and 0.176, and the larger part of that gain came from the
recalibration rather than from the retraining.

\textbf{Scientific Contribution} This work separates the temporal validation penalty into
discrimination and calibration and shows the calibration component survives
matching on both base rates, so it does not reduce to prevalence shift. It also does not reduce to the amount of covariate shift, since splits built to
match the measured shift remain less affected than the time split. Periodic retraining with recalibration is the one remedy that closes most of
the gap, which turns a reporting recommendation into an operating procedure.}

%%================================%%
%% Sample for structured abstract %%
%%================================%%

% \abstract{\textbf{Purpose:} The abstract serves both as a general introduction to the topic and as a brief, non-technical summary of the main results and their implications. The abstract must not include subheadings (unless expressly permitted in the journal's Instructions to Authors), equations or citations. As a guide the abstract should not exceed 200 words. Most journals do not set a hard limit however authors are advised to check the author instructions for the journal they are submitting to.
% 
% \textbf{Methods:} The abstract serves both as a general introduction to the topic and as a brief, non-technical summary of the main results and their implications. The abstract must not include subheadings (unless expressly permitted in the journal's Instructions to Authors), equations or citations. As a guide the abstract should not exceed 200 words. Most journals do not set a hard limit however authors are advised to check the author instructions for the journal they are submitting to.
% 
% \textbf{Results:} The abstract serves both as a general introduction to the topic and as a brief, non-technical summary of the main results and their implications. The abstract must not include subheadings (unless expressly permitted in the journal's Instructions to Authors), equations or citations. As a guide the abstract should not exceed 200 words. Most journals do not set a hard limit however authors are advised to check the author instructions for the journal they are submitting to.
% 
% \textbf{Conclusion:} The abstract serves both as a general introduction to the topic and as a brief, non-technical summary of the main results and their implications. The abstract must not include subheadings (unless expressly permitted in the journal's Instructions to Authors), equations or citations. As a guide the abstract should not exceed 200 words. Most journals do not set a hard limit however authors are advised to check the author instructions for the journal they are submitting to.}

%% [TODO] 저널 지침 확인 후 확정.
\keywords{calibration, temporal validation, drug-induced liver injury,
drug-induced cardiotoxicity, conformal prediction, model evaluation}

%%\pacs[JEL Classification]{D8, H51}

%%\pacs[MSC Classification]{35A01, 65L10, 65L12, 65L20, 65L70}

\maketitle
